\section{固定引戻し量子圏のhull：再利用可能な論証}
\label{notes:pullback-fixed-hulls-v37}

\subsection{確定した対象と結論}

$m\ge3$、$N_m=\binom{m+2}{2}$、$c=xy(x+y)$、$L=Ax+By$ とする。
$\deg(x,y,z,w)=(1,1,1,m)$ に対し
\[
 \delta(x)=\frac{c_y}{3}+\frac{m}{3(m+3)}xL,\quad
 \delta(y)=-\frac{c_x}{3}+\frac{m}{3(m+3)}yL,\quad
 \delta(w)=u x^{m+1}-\frac{m}{m+3}Lw
\]
と置き、$S=\mathbf C[x,y,w][z;s\delta]$ の標準 $m+3$ 対象列の有限代数を
$\mathcal A$ とする。$u,s\ne0$ とし、下記の明示開条件を課すと、
\[
 R_{\mathcal A}\simeq
 \mathbf C[[t_1,t_2,a,b]]/(a^{m+2},ab,b^2).
\]
二つの滑らかな変数は実際のOre係数 $A,B$ の変化によって実現される。
$HH^\bullet$ の次元は $m=3$ で $(1,2,4,9)$、$m\ge4$ で $(1,2,4,6)$、
次数三を超える群は零である。
これは全結合変形のhullであり、次数付き座標環の小さなパラメータ族だけの計算ではない。
ただし対象は標準有限代数であって、無限次元環 $S$ の全非次数付き変形とは区別する。

十分な開集合は次の通りである。
\[
 \kappa=\left(1+\frac m3(1-B),\,1+\frac m3(1+A),\,
                  1+\frac m3(1-A+B)\right),
\]
\[
 (B-2A)(A-2B)(A+B)\kappa_1\kappa_2\kappa_3
 [m(B+2)+6]\prod_{k=0}^{m}[m(B+2)-3k]\ne0,
\]
かつ $\kappa$ が和 $m+3$ の正整数三つからなる有限集合を避けることを要求する。
$(A,B)=(2,5)$ は全ての $m\ge3$ でこれを満たす。
この条件集合の外における例外パラメータは分類していない。

\subsection{非零混合点への局所環の移動}

共点三直線での既存の全Poisson局所環は
\[
 \mathbf C[[s_1,\ldots,s_{N_m+8},a,b,u]]/(ab,b^2,a^{m+2}u)
\]
である。$u=0$ での形式級数に $u_0\ne0$ を直接代入することはできない。
再利用すべき性質は、証明全体が混合変数 $F$ に関して線形である点である。
$(c,D)$ の切片は $F$ に依存せず、Schur消去、移動像の直和因子の分離、
残余座標 $u$ の選択は $F$ の係数について線形である。
従って基底 $(c,D)$ だけを完備化し、$F$ は多項式変数のままとして消去できる。
その後 $F=u_0x^{m+1}\partial_z$ で完備化し、$u=u_0+v$ と置くと
\[
 \widehat{\mathcal O}_{\mathscr P_m,\pi}
 =T_m[[z_1,\ldots,z_{N_m+9}]],\qquad
 T_m=\mathbf C[a,b]/(a^{m+2},ab,b^2)
\]
を得る。これが非零混合点での定理の入力となる。

開条件の由来も明示できる。二変数移動の双対角行列は
\[
 V(x^{m-k}y^k)=
 \left(\frac{m(A+1)}3-k\right)x^{m+1-k}y^k+
 \left(\frac{m(B+2)}3-k\right)x^{m-k}y^{k+1}
\]
であり、第二係数の全てが非零なら単射かつ余核は $x^{m+1}$ で生成される。
最高重量 $m+1$ で、源を $z^{m+1}\partial_x,z^{m+1}\partial_y$、
消す標的を $xyz^{m+1},y^2z^{m+1}$ にすると、小行列式は
\[
 -[m(1-B)+3][m(B+2)+6]/9
\]
となる。低い重量の標的全射性は $\kappa_i\ne0$ から分かる。
従って不可視の追加小行列条件を置く必要はない。

\subsection{二つの安定化方向と二つの変形方向}

重み付き座標変換のLie代数は
$\mathfrak{sl}_3\oplus\mathbf Cw\partial_w\oplus S_m\partial_w$
で、次元 $N_m+9$ である。
$c$ の軌道が六次元、移動像が $N_m$ 次元、非零混合係数の倍率が一方向なので、
全Poisson軌道次元は $N_m+7$ である。
安定化Lie代数は $x\partial_z,y\partial_z$ の二次元である。

量子側でも $z\mapsto z+tx$、$z\mapsto z+ty$ は正確な自己同型である。
この二つの導分の生成元上の値は $z\mapsto x,y$ であり、他は零である。
正規順序のPBW多項式上で通常の $x\partial_z,y\partial_z$ と同一視してはならない。
高い $z$ 冪にはOre補正が生じる。
標準有限代数の冪等元を固定する内的導分は対角スカラーから来るので、
この二方向は独立な外的導分となる。
Hochschild複体の最初の差分 $[\pi,-]$ による上界と合わせて $HH^1=2$ を得る。

可換点の基底変換軌道次元を $b_0$ とすると、量子点での値は
$b_\eta=b_0+N_m+7$ である。
全結合積スキームのgeneric局所環には $b_0+N_m+9$ 個の滑らかな変数があるため、
二つがhullに残る。Artin環だけを最終hullとする結論は誤りである。

さらに $A,B$ の二つのKodaira--Spencer類は独立である。
Poisson側では、$c$ を固定する $\mathfrak{sl}_3$ 導分は二つの $z$ 移動だけで、
それらも $w$ 変換も $D$ を変えない。一方 $\partial_A D,\partial_B D$ は独立である。
量子側で可縮部分を除いた複体を取ると、次数一の微分は
$d_h=h(D_1+O(h))$ であり、$HH^1$ が先頭差分とgeneric fiberの双方で二次元なので、
$d_h/h$ と $D_1$ の階数が一致する。
従って $\operatorname{im}(d_h/h)$ は $\mathbf C[[h]]$ 上の直和因子である。
$A,B$ のKSコサイクルを $h$ で割ってから特殊化すると $\partial_A\pi,\partial_B\pi$ となる。
直和因子性によってgeneric fiberでの従属は先頭項の従属を引き起こすので、独立性を得る。
この議論は、最初のコホモロジー差分で階数が確定した明示量子族へ一般に再利用できる。

\subsection{形式移送で必要な三つの一般原理}

第一に、斉次錐の点 $\pi$ の局所環が $T[[z_1,\ldots,z_d]]$、$T$ が有限局所環なら、
$p(h)=h\pi+O(h^2)$ のgeneric局所環も $T_{\mathbf C((h))}[[z_1,\ldots,z_d]]$ である。
$p/h$ を $\pi$ で再中心化し、形式級数を $h$ 進に評価する。
被約環 $\mathbf C[[h]]$ への弧は $T$ の冪零元を殺すため、動くのは滑らかな座標だけである。

第二に、$T[[z_1,\ldots,z_d]]$ 上の滑らかな群作用で閉点軌道次元が $q$ なら、
横断切片を $T[[z_{q+1},\ldots,z_d]]$ に取れる。
被約な軌道上では $T$ の冪零元が消えるので、軌道接空間は $z$ 方向だけに入る。
$q$ 個の群座標を選び、最初の $q$ 個の $z$ を零にする方程式へ形式逆関数定理を使う。
これはArtin厚みの同型型を保存したまま軌道を除く論証であり、
滑らかな変数の数だけを数えて形式的な消去を推測する方法とは区別される。

第三に、複素係数の有限埋込次元完備Noether局所環二つが、拡大体上で形式同型なら、
元の複素係数でも非標準的に同型である。
定義関係、形式同型、逆同型、および関係所属の証人に現れる係数は可算個である。
定義関係の係数を含む可算部分体 $k_0\subset\mathbf C$ を固定し、
残りの係数を添加した可算体を $k_0$ 上で $\mathbf C$ に埋め込む。
$\mathbf C$ の超越次数と代数的閉性がその埋込みを保証する。
全係数恒等式が保たれるため、互いに逆な複素係数の形式写像を得る。
これは $h=1$ を形式級数へ代入する操作ではない。

\subsection{再利用する環の不変量と未決の範囲}

$T_m$ の基底は $1,a,\ldots,a^{m+1},b$、socleは
$\langle a^{m+1},b\rangle$ である。
hullは二次元Cohen--Macaulay、型二、重複度 $m+3$、埋込次元四であり、
非Gorensteinかつ非完全交叉である。
障害の三つの最小関係は $a^{m+2},ab,b^2$ である。
従って純粋な $a$ 方向は次数 $m+2$ で初めて止まり、$b$ は二次で止まる。

残る有用な問題は、$A,B$ 間の大域的な同値分類、明示開集合の外側のhull、
および二次元安定化群の非被約hull上の作用である。
本結果はmarked miniversal環の同型型を決定するが、
自己同型を含む変形groupoidや、大域的なcoarse moduliを分類しない。

検証コードは標準有限代数の矢印像を未知数として、全最小Ore関係のLeibniz式を有理数上で解く。
$m=3,4,5,6$ の $HH^1=2$、二つの明示外的導分、$A,B$ のKS類の独立性、
および別構成のPoisson Lie微分行列の軌道階数を照合する。
これらの有限計算は全重みの構造証明の監査であり、外挿の根拠には使わない。
